\textbf{Актуальность темы исследования}

В задачах численного моделирования, анализа динамических систем и визуализации векторных полей часто требуется многократно вычислять значения сложных отображений и решений дифференциальных уравнений на некоторой области параметров. Если исходная функция задана аналитически, но её вычисление дорого, либо определяется как результат численного интегрирования, прямой пересчёт в каждой точке становится узким местом интерактивных и серверных систем. Особенно остро эта проблема проявляется при построении плотных карт фазовых траекторий, визуализации эволюции двумерных систем и обработке большого числа пользовательских запросов в режиме реального времени.

Одним из эффективных способов снижения вычислительных затрат является аппроксимация отображения на основе разреженных сеток. Такой подход позволяет существенно уменьшить число узлов по сравнению с полными тензорными сетками и при этом сохранить приемлемую точность аппроксимации. Дополнительный интерес представляет адаптивное уточнение сетки, при котором вычислительные ресурсы концентрируются только в тех областях, где функция имеет сложное поведение. В сочетании с параллельным построением и быстрым рекурсивным вычислением значений это делает разреженные сетки перспективным инструментом для высокопроизводительных вычислительных сервисов.

\textbf{Степень разработанности проблемы}

Теория разреженных сеток и иерархических базисов хорошо развита в численном анализе и применяется для решения задач интерполяции, квадратур и аппроксимации многомерных функций. В классических работах показано, что использование иерархических представлений позволяет бороться с «проклятием размерности» эффективнее, чем прямое построение полных сеток. Для задач адаптивной интерполяции широко используются линейные базисные функции, так как они просты в реализации и анализе. Однако квадратичные базисы позволяют точнее описывать гладкие зависимости при сопоставимом числе уровней детализации, что делает их особенно интересными для практических приложений.

В существующих программных реализациях упор обычно делается либо на чисто математические пакеты, либо на узкоспециализированные численные библиотеки без развитого прикладного интерфейса для удалённого вызова и визуального анализа. Кроме того, далеко не все решения сочетают одновременно адаптивное уточнение сетки, композиционное применение интерполянтов, параллельное построение узлов, поддержку дифференциальных уравнений и удобную визуализацию результата. Это определяет необходимость разработки целостного программного комплекса, объединяющего математическое ядро, сервер вычислений и пользовательский интерфейс.

\textbf{Проблема исследования}

При численном анализе сложных отображений возникает противоречие между точностью и стоимостью вычислений. С одной стороны, повышение точности требует увеличения количества узлов аппроксимации и числа обращений к исходной функции. С другой стороны, при интерактивной работе пользователя или обслуживании удалённых запросов время ответа должно оставаться ограниченным. Использование полных сеток быстро становится непрактичным уже в двумерных и трёхмерных задачах, а последовательное построение адаптивной структуры данных не позволяет полностью задействовать возможности современных многоядерных систем. Следовательно, требуется метод и программная реализация, обеспечивающие баланс между точностью, скоростью построения, масштабируемостью и удобством интеграции в прикладные системы.

\textbf{Цель работы}

Целью выпускной квалификационной работы является разработка и исследование программного комплекса для интерполяции отображений и решений дифференциальных уравнений на основе адаптивных разреженных сеток с поддержкой квадратичного базиса и параллельных алгоритмов построения.

\textbf{Задачи исследования}

Для достижения поставленной цели в работе были сформулированы и решены следующие задачи:

\begin{enumerate}
    \item Исследовать математические основы иерархической интерполяции на разреженных сетках, включая линейный и квадратичный базисы.
    \item Разработать структуру данных для представления адаптивной разреженной сетки, узлов сетки и входных точек различных размерностей.
    \item Реализовать алгоритм построения адаптивной сетки с использованием локального критерия уточнения по величине иерархических коэффициентов и контрольных точек.
    \item Реализовать последовательный и параллельный режимы построения сетки и сравнить их на наборе тестовых задач.
    \item Обеспечить возможность аппроксимации как обычных отображений, так и операторов перехода, возникающих при численном решении систем дифференциальных уравнений.
    \item Разработать серверный и клиентский уровни программного комплекса для удалённого вычисления и визуализации результатов.
    \item Выполнить тестирование, бенчмаркинг и анализ эффективности разработанного решения.
\end{enumerate}

\textbf{Объект и предмет исследования}

\textbf{Объектом исследования} являются методы численной интерполяции многомерных функций и решений дифференциальных уравнений на основе разреженных сеток. \\
\textbf{Предметом исследования} являются параллельные алгоритмы построения адаптивных разреженных сеток, структуры данных вычислительного ядра и программные средства визуализации результатов интерполяции.

\textbf{Методы исследования}

В работе использовались методы численного анализа, теории интерполяции, математического моделирования динамических систем, а также методы проектирования программного обеспечения. Для решения задач аппроксимации применялись иерархические базисные функции, рекурсивные алгоритмы обхода структуры сетки и численное интегрирование обыкновенных дифференциальных уравнений методом Рунге--Кутты--Фельберга. Для оценки качества реализации использовались модульное тестирование, тестирование на характерных функциях и бенчмаркинг.

\textbf{Научная новизна и практическая значимость}

Научная новизна работы состоит в объединении адаптивного подхода к построению разреженных сеток с параллельным режимом генерации узлов и использованием квадратичного базиса в прикладной системе, ориентированной на интерполяцию отображений и решений дифференциальных уравнений. Предложенная реализация использует общий вычислительный механизм как для прямой аппроксимации функции, так и для итеративного построения композиции отображений, что позволяет применять один и тот же аппарат к широкому классу задач.

Практическая значимость работы заключается в создании законченного программного комплекса, включающего вычислительное ядро на Go, gRPC-сервис для удалённых вызовов и клиентское приложение на Flutter для визуального анализа построенной сетки. Разработанное решение может быть использовано для исследования двумерных динамических систем, интерактивной визуализации векторных полей, построения приближённых операторов эволюции и демонстрации методов адаптивной интерполяции в учебных и исследовательских целях.

\textbf{Апробация результатов}

Результаты работы не проходили апробацию в форме публикаций и докладов на конференциях.